% Scheletro di lavoro, non testo definitivo.
% Fonti e passi sono proposte di lettura: verificarli sul testo greco prima della stesura.
% Seguira il metodo: documento originale, traduzione, esame dell'argomentazione, ricostruzione critica.
\chapter{L'anima, l'immortalità e l'eros}
\label{cap:platone:anima-eros}

% Problema: Come viene definito il rapporto fra anima, corpo e filosofia?
% Fonti da verificare: Fedone 64c–69e; Fedro 245c–246a.
\section{La natura dell'anima}

\subsection{L'anima e il corpo nel Fedone}

\subsection{L'anima come principio di movimento nel Fedro}

% Problema: Quale forza dimostrativa hanno gli argomenti sull'immortalità?
% Fonti da verificare: Fedone 70c–84b, 102a–107b.
\section{L'immortalità dell'anima}

\subsection{Gli argomenti del Fedone}

\subsection{Obiezioni, risposte e limiti}

% Problema: Come spiegare il conflitto fra desideri e ragione?
% Fonti da verificare: Repubblica IV 435b–441c; Fedro 246a–249d.
\section{La struttura dell'anima}

\subsection{Le tre componenti nella Repubblica}

\subsection{L'immagine del carro alato}

% Problema: In che modo il desiderio può trasformarsi in ricerca del vero?
% Fonti da verificare: Simposio 201d–212b; Fedro 249d–257b.
\section{L'eros e l'ascesa filosofica}

\subsection{Diotima e l'ascesa al bello}

\subsection{Eros, memoria e contemplazione}
